\section{Discussion}
\label{sec:discussion}

\paragraph{What F(AI)\textsuperscript{2}R is not.}
It is not a vendor stack: substituting Turtle for JSON-LD, GitHub
Actions for any other CI, or any one model for any other, leaves the
eight practices intact. It is not a guarantee: a team that adopts
the pattern can still ship a bad paper. The pattern reduces specific
failure modes to specific remediations and surfaces residual ones; it
does not promise that what survives is true.

\paragraph{Limitations of a single worked example.}
The pattern has been exercised on two papers (the source repository
and this one), authored by the same researcher. The sample is too
small to support claims about the pattern's behaviour at scale, in
adversarial review, or in collaborations whose contributors do not
share the same toolchain. We expect verification overhead to be
sensitive to the maturity of the harness; we do not have data to
quantify the dependence.

\paragraph{Submission conditions.}
We will submit F(AI)\textsuperscript{2}R to the Research Data
Alliance and to FORCE11 as a candidate community recommendation when
two preconditions are met: (i) human-author consent for
community-recommendation submission is recorded in
\texttt{doc/publication-consent.md}; and (ii) at least one external
case study has adopted the pattern on a paper unrelated to either
the source repository or this one, and the resulting artefacts are
public under a permissive licence. The conditions are the minimum we
need to distinguish a single-author idiosyncrasy from a transferable
pattern.

\paragraph{Dual use.}
A pipeline of prompts, a finite-state ladder, and a structured
provenance schema are an instruction set adoptable by bad-faith
authors as readily as by good-faith ones.
F(AI)\textsuperscript{2}R \emph{makes the dual-use surface visible
in a form that can be reasoned about}; it does not eliminate it.
The mitigation is community: a bad-faith adoption is detectable by
inspection of the artefacts the pattern requires.

\paragraph{Naming the practice.}
\emph{Naming what we are already doing is the first step toward
surrendering it to peer scrutiny.} The eight practices were not
invented for this paper; they are practices the source repository's
author had gradually adopted by attrition, after each
individually-unoriginal practice paid for itself against a specific
failure mode. The act of naming, here, makes the practices
contestable in the open literature and lets a third party adopt
them without re-deriving each from the failure modes that motivated
it.
